% ============================================================
% sec/90meetings.tex - apendice: reuniones y sesiones de trabajo
% ============================================================
\section{Reuniones y sesiones de trabajo}
\label{sec:meetings}

\subsection{\'Indice}

Las reuniones de la Etapa 1 se hicieron por llamada en Discord y quedaron
registradas con capturas. Desde la Etapa 2 la coordinaci\'on pas\'o a
sesiones de trabajo por rama y a acuerdos escritos en
\code{docx/\allowbreak{}03Evidencia/\allowbreak{}}, que funcionan como actas. La duraci\'on de las
llamadas es la que muestra el contador de la captura, as\'i que es una cota
inferior.

\begin{table}[H]
  \centering
  \caption{Reuniones y sesiones registradas.}
  \label{tab:meetingindex}
  \footnotesize
  \begin{tabularx}{\textwidth}{@{}l P{2.0cm} P{3.4cm} Y@{}}
    \toprule
    \textbf{Fecha} & \textbf{Duraci\'on} & \textbf{Participantes} & \textbf{Tema principal}\\
    \midrule
    21/08 & 20 min o m\'as & Adriel, Daniel, Sebasti\'an & Primera reuni\'on; conjuntos de datos en Kaggle\\
    26/08 & 15 min o m\'as & Adriel, Sebasti\'an & Confirmaci\'on del tema\\
    30/08 & 13 min o m\'as & Adriel, Daniel, Sebasti\'an & Cierre de la confirmaci\'on y reparto de ejes\\
    06/09 & 1~h 38~min o m\'as & Adriel, Daniel, Sebasti\'an & Revisi\'on final del art\'iculo de la Etapa 1\\
    12/10 & Acta escrita & Daniel & Acuerdos t\'ecnicos de modelado\\
    28/10 & Acta escrita & Daniel & Selecci\'on de B3, definici\'on de H2 y protocolo de LightGBM\\
    02/11 & Acta escrita & Sebasti\'an & Protocolo de evaluaci\'on final\\
    \bottomrule
  \end{tabularx}
\end{table}

\subsection{Actas}

\jentry{2026-08-21}{Primera reuni\'on del equipo}{Adriel, Daniel, Sebasti\'an}{---}

\noindent\textbf{Temas.}
\begin{itemize}
  \item Revisi\'on de conjuntos de datos p\'ublicos en Kaggle.
  \item Organizaci\'on del repositorio.
\end{itemize}

\noindent\textbf{Acuerdos.}
\begin{itemize}
  \item Adriel crea el repositorio con las reglas de ramas y commits
        (cumplido el mismo d\'ia, PR \#1).
  \item El equipo considera las repeticiones del PTCG AI Battle Challenge
        como corpus candidato.
\end{itemize}

\jentry{2026-08-26}{Confirmaci\'on del tema}{Adriel, Sebasti\'an}{---}

\noindent\textbf{Acuerdos.}
\begin{itemize}
  \item Track~A: probabilidad de victoria en el Pok\'emon TCG desde el estado
        de la partida.
  \item Se deja la direcci\'on de aprendizaje por refuerzo explorada
        antes.
\end{itemize}

\jentry{2026-08-30}{Cierre de la confirmaci\'on del \textit{track}}%
       {Adriel, Daniel, Sebasti\'an}{---}

\noindent\textbf{Acuerdos.}
\begin{itemize}
  \item Base de investigaci\'on de la Etapa 1 cerrada (PR \#2).
  \item Revisi\'on de literatura repartida por ejes: Daniel el A, Adriel el
        B y Sebasti\'an el C.
\end{itemize}

\jentry{2026-09-06}{Revisi\'on final del art\'iculo}{Adriel, Daniel, Sebasti\'an}{---}

\noindent\textbf{Temas.}
\begin{itemize}
  \item Lectura completa del art\'iculo antes de la entrega.
  \item Ajustes de formato, de la tabla comparativa y de la metodolog\'ia.
\end{itemize}

\noindent\textbf{Acuerdos.}
\begin{itemize}
  \item Integrar \code{02PaperAndDelivery} y \code{Stage1} a \code{master}
        (PR \#4 y \#5) y entregar.
\end{itemize}

\jentry{2026-10-12}{Acuerdos t\'ecnicos de modelado}{Daniel}{04MainAIModelTrainingAndDevelopment}

Acta: \code{12-\allowbreak{}10-\allowbreak{}26\_\allowbreak{}AcuerdosTecnicosModelado.md}. Fij\'o la f\'ormula de
B1, la validaci\'on cruzada agrupada, los pesos por partida, la b\'usqueda de
B2 en cuadr\'icula, las b\'usquedas de B3 y B4 con Optuna y la parada
temprana por partidas. Dej\'o abiertas las semillas m\'ultiples y el tipo de
intervalo, que se cerraron en el protocolo final (una semilla; percentil del
95\,\% con 1000 remuestreos).

\jentry{2026-10-28}{B3, H2 y LightGBM}{Daniel}{04MainAIModelTrainingAndDevelopment}

Actas: \code{28-\allowbreak{}10-\allowbreak{}26\_\allowbreak{}SeleccionB3yDefinicionH2.md} y
\code{28-\allowbreak{}10-\allowbreak{}26\_\allowbreak{}ProtocoloLightGBM.md}. Fijaron el macro-AUC para elegir B3,
la definici\'on de las cuatro entradas de H2 y la regla para reemplazar B4
por LightGBM, todo antes de ejecutar.

\jentry{2026-11-02}{Protocolo de evaluaci\'on final}{Sebasti\'an}{05FinalModelChecks}

Acta: \code{02-\allowbreak{}11-\allowbreak{}26\_\allowbreak{}ProtocoloEvaluacionFinal.md}. Fij\'o el orden de las
barreras, las reglas de decisi\'on de H1, H2 y H3, la calibraci\'on decidida
con \code{calib} y los umbrales de la revisi\'on de fuga. Se congel\'o con un
commit antes de abrir prueba.